\chapter{Thực nghiệm}
\label{chap:experiments}

Chương này trình bày thiết lập thực nghiệm và kết quả đánh giá hệ thống Self-Evolving AI Agents trên benchmark NIKA. Mục tiêu chính là so sánh agent ReAct thông thường với agent ReAct được bổ sung Procedural Memory và Tool Refinement.

\section{Thiết lập thực nghiệm}
\label{sec:experiment_setup}

Thực nghiệm so sánh hai cấu hình agent:
\begin{itemize}
    \item \textbf{Baseline Agent:} ReAct agent tiêu chuẩn, không sử dụng bộ nhớ tiến hóa và không tinh chỉnh tài liệu công cụ.
    \item \textbf{Evolved Agent:} ReAct agent tích hợp Procedural Memory và Tool Refinement, cho phép tái sử dụng kỹ năng chẩn đoán và tài liệu công cụ đã được cải thiện từ các episode trước.
\end{itemize}

Cả hai cấu hình sử dụng cùng mô hình nền \texttt{openai/gpt-oss-120b}, cùng tập MCP servers và giới hạn \texttt{max\_steps=20}. Với Baseline Agent, mỗi incident được xử lý độc lập. Với Evolved Agent, kinh nghiệm từ các incident trước được cập nhật ngoại tuyến và sử dụng cho các incident sau.

Quy trình chạy benchmark gồm bốn bước: khởi động lab Kathará, tiêm lỗi xác định, chạy agent chẩn đoán và đánh giá kết quả. Môi trường thực nghiệm sử dụng Python 3.12, LangGraph/LangChain cho orchestration, FastMCP cho MCP servers và Kathará cho ảo hóa mạng.

\section{Benchmark và chỉ số đánh giá}
\label{sec:benchmark_dataset}

Hai lần chạy benchmark gồm 125 case chung thuộc 6 nhóm nguyên nhân gốc rễ trên 8 kịch bản mạng: OSPF enterprise, BGP CLOS data center, SDN CLOS, RIP/VPN và P4/BMv2. Các nhóm lỗi bao gồm link failures, end-host failures, misconfigurations, resource contention, network node errors và network under attack; ngoài ra có các case \texttt{no\_fault} để kiểm tra khả năng không báo động sai.

Mỗi incident có ground truth gồm ba trường: trạng thái bất thường \texttt{is\_anomaly}, danh sách thiết bị lỗi \texttt{faulty\_devices} và tên nguyên nhân gốc rễ \texttt{root\_cause\_name}. Bảng kết quả tổng thể được tính trực tiếp từ hai batch \texttt{benchmark\_evaluate-0001} và \texttt{benchmark\_evaluate-0002}.

Hệ thống được đánh giá theo ba nhóm chỉ số:
\begin{itemize}
    \item \textbf{Detection:} Agent xác định đúng mạng có sự cố hay không.
    \item \textbf{Localization:} Agent xác định đúng thiết bị lỗi, đo bằng accuracy và F1 trên tập thiết bị.
    \item \textbf{RCA:} Agent xác định đúng nguyên nhân gốc rễ, đo bằng accuracy và F1.
\end{itemize}

Một incident được xem là thành công toàn diện khi agent đồng thời phát hiện đúng sự cố, định vị đúng toàn bộ thiết bị lỗi và xác định đúng nguyên nhân gốc rễ. Ngoài độ chính xác, thực nghiệm cũng đo số bước suy luận, số lần gọi công cụ và lượng token trung bình.

\section{Kết quả thực nghiệm}
\label{sec:results}

Bảng \ref{tab:main_results} trình bày kết quả chung của cấu hình ReAct + GPT-OSS-120B, so sánh Baseline Agent và Evolved Agent.

\begin{table}[htbp]
\centering
\caption{Kết quả chung của ReAct + GPT-OSS-120B}
\label{tab:main_results}
\begin{tabular}{|L{5cm}|C{3cm}|C{3cm}|}
\hline
\textbf{Chỉ số} & \textbf{Baseline} & \textbf{Evolved} \\
\hline
Detection Score & 0.840 & 0.872 \\
\hline
Localization Accuracy & 0.368 & 0.384 \\
\hline
Localization F1 & 0.411 & 0.459 \\
\hline
RCA Accuracy & 0.304 & 0.376 \\
\hline
RCA F1 & 0.304 & 0.387 \\
\hline
Incident Success Rate & 26/125 (20.8\%) & 35/125 (28.0\%) \\
\hline
Steps trung bình & 15.29 & 12.74 \\
\hline
Tool calls trung bình & 13.60 & 11.09 \\
\hline
Tool error rate & 11/1700 (0.65\%) & 0/1386 (0.00\%) \\
\hline
Input tokens (tổng) & 15.944M & 17.992M \\
\hline
Output tokens (tổng) & 0.697M & 0.708M \\
\hline
Total tokens (tổng) & 16.641M & 18.699M \\
\hline
Tổng thời gian chạy & 207.7 phút & 222.6 phút \\
\hline
\end{tabular}
\end{table}

Kết quả cho thấy Evolved Agent cải thiện trên hầu hết các chỉ số chẩn đoán. Incident success rate tăng từ 20.8\% lên 28.0\% (26/125 lên 35/125), tương đương tăng 34.6\%. RCA F1 tăng từ 0.304 lên 0.387, Localization F1 tăng từ 0.411 lên 0.459. Đồng thời, agent tự tiến hóa dùng ít bước hơn (giảm 16.7\%), gọi công cụ ít hơn (giảm 18.5\%) và giảm tool error rate từ 0.65\% xuống 0.00\%.

Mức cải thiện rõ nhất xuất hiện ở các lỗi cấu hình host như thiếu IP, sai IP hoặc sai gateway. Với các lỗi này, Procedural Memory giúp agent tái sử dụng quy trình kiểm tra cấu hình host, còn Tool Refinement giúp mô tả công cụ rõ hơn về tham số và cách diễn giải output. Ở các lỗi BGP/OSPF phức tạp, SDN/P4 hoặc resource contention, cải thiện vẫn còn hạn chế vì cần suy luận nhiều tầng và phụ thuộc mạnh vào bằng chứng từ nhiều công cụ.

\subsection{Các bảng và hình bổ sung dự kiến}
\label{sec:planned_result_artifacts}

Các bảng và hình trong phần này là placeholder để cố định bố cục báo cáo. Các giá trị được đánh dấu \textbf{minh họa} và sẽ được thay bằng kết quả thực nghiệm chính thức sau khi hoàn tất các batch benchmark tương ứng.

\begin{figure}[htbp]
\centering
\fbox{%
\begin{minipage}{0.92\linewidth}
\centering
\small
\textbf{Minh họa xu hướng metric theo số case tăng dần}\\[0.4em]
\begin{tabular}{c|cc|cc}
\hline
\textbf{Case} & \multicolumn{2}{c|}{\textbf{RCA F1}} & \multicolumn{2}{c}{\textbf{Localization F1}} \\
 & Baseline & Evolved & Baseline & Evolved \\
\hline
25 & 0.22 & 0.43 & 0.34 & 0.55 \\
50 & 0.28 & 0.42 & 0.39 & 0.53 \\
75 & 0.30 & 0.40 & 0.40 & 0.49 \\
100 & 0.31 & 0.39 & 0.41 & 0.47 \\
125 & 0.30 & 0.39 & 0.41 & 0.46 \\
\hline
\end{tabular}\\[0.4em]
\footnotesize Đường cong dự kiến thể hiện cải thiện nhanh ở nhóm case đầu do có nhiều case liên quan, sau đó chậm lại khi gặp nhóm case xa hơn.
\end{minipage}}
\caption{Placeholder figure: xu hướng RCA và Localization theo số case tăng dần. Số liệu minh họa, chưa dùng làm kết quả chính thức.}
\label{fig:cumulative_metrics_placeholder}
\end{figure}

\begin{figure}[htbp]
\centering
\fbox{%
\begin{minipage}{0.92\linewidth}
\centering
\small
\textbf{Minh họa độ ổn định qua 3 lần chạy ReAct + GPT-OSS-120B}\\[0.4em]
\begin{tabular}{c|cc|cc}
\hline
\textbf{Run} & \multicolumn{2}{c|}{\textbf{Incident Success}} & \multicolumn{2}{c}{\textbf{RCA F1}} \\
 & Baseline & Evolved & Baseline & Evolved \\
\hline
1 & 20.8\% & 28.0\% & 0.304 & 0.387 \\
2 & 19.6\% & 27.2\% & 0.296 & 0.381 \\
3 & 21.6\% & 29.6\% & 0.311 & 0.395 \\
\hline
\end{tabular}\\[0.4em]
\footnotesize Figure chính thức sẽ dùng error bar hoặc box plot để thể hiện trung bình và độ lệch chuẩn giữa các lần chạy.
\end{minipage}}
\caption{Placeholder figure: so sánh độ ổn định giữa baseline và evolved qua 3 lần chạy. Số liệu minh họa, chưa dùng làm kết quả chính thức.}
\label{fig:stability_placeholder}
\end{figure}

\begin{table}[htbp]
\centering
\caption{Placeholder table: so sánh ReAct, Plan\&Execute và Reflexion trên GPT-OSS-120B. Số liệu minh họa, chưa dùng làm kết quả chính thức.}
\label{tab:agent_strategy_placeholder}
\begin{tabular}{|L{3.2cm}|C{2.2cm}|C{2.2cm}|C{2.2cm}|}
\hline
\textbf{Chiến lược agent} & \textbf{Incident Success} & \textbf{RCA F1} & \textbf{Tool calls} \\
\hline
ReAct & 28.0\% & 0.387 & 11.09 \\
\hline
Plan\&Execute & 25.6\% & 0.361 & 12.74 \\
\hline
Reflexion & 26.4\% & 0.374 & 13.28 \\
\hline
\end{tabular}
\end{table}

\begin{figure}[htbp]
\centering
\fbox{%
\begin{minipage}{0.92\linewidth}
\centering
\small
\textbf{Minh họa so sánh mô hình nền với ReAct}\\[0.4em]
\begin{tabular}{c|cc|cc}
\hline
\textbf{Model} & \multicolumn{2}{c|}{\textbf{Baseline}} & \multicolumn{2}{c}{\textbf{Evolved}} \\
 & Success & RCA F1 & Success & RCA F1 \\
\hline
GPT-OSS-120B & 20.8\% & 0.304 & 28.0\% & 0.387 \\
Qwen3.5-122B-A10B-FP8 & 18.4\% & 0.286 & 25.6\% & 0.365 \\
\hline
\end{tabular}\\[0.4em]
\footnotesize Figure chính thức sẽ dùng grouped bar chart để so sánh ReAct trên hai model nền.
\end{minipage}}
\caption{Placeholder figure: so sánh ReAct trên GPT-OSS-120B và Qwen3.5-122B-A10B-FP8. Số liệu minh họa, chưa dùng làm kết quả chính thức.}
\label{fig:model_comparison_placeholder}
\end{figure}

\section{Phân tích hiệu quả và hạn chế}
\label{sec:discussion}

Kết quả thực nghiệm xác nhận rằng hướng tiếp cận tự tiến hóa có thể cải thiện hiệu quả chẩn đoán mà không cần fine-tuning mô hình nền. Procedural Memory đóng vai trò tích lũy quy trình chẩn đoán ở mức chiến lược, giúp agent chọn thứ tự kiểm tra hợp lý hơn. Tool Refinement cải thiện hiểu biết ở mức công cụ, giảm lỗi gọi tham số và giúp agent diễn giải output chính xác hơn.

Đổi lại, Evolved Agent tiêu thụ nhiều token hơn Baseline Agent: tổng token tăng từ 16.641M lên 18.699M (+12.4\%), trong đó input tokens tăng từ 15.944M lên 17.992M và output tokens tăng nhẹ từ 0.697M lên 0.708M. Tổng thời gian chạy tăng từ 207.7 lên 222.6 phút (+7.2\%). Đây là chi phí chấp nhận được trong bối cảnh incident success rate tăng 34.6\%, nhưng vẫn cần tối ưu nếu triển khai trong môi trường vận hành quy mô lớn.

Một hạn chế khác là khả năng tổng quát hóa chưa đồng đều giữa các nhóm sự cố. Các kỹ năng học từ incident trước có thể phù hợp tốt với lỗi host hoặc lỗi cấu hình đơn giản, nhưng chưa đủ mạnh cho các sự cố liên quan đến định tuyến động, tấn công mạng hoặc tương tác phức tạp giữa nhiều thành phần. Điều này cho thấy cần tiếp tục cải thiện cơ chế trừu tượng hóa kỹ năng và chọn lọc kinh nghiệm trong các phiên bản tiếp theo.
